% SimulationStyle.tex
% Academic Beamer style for the 20260918 Science Quiz Simulation.

\usetheme{Madrid}
\useinnertheme{rounded}
\usefonttheme{professionalfonts}

\usepackage{kotex}
\usepackage[scaled=0.94]{helvet}
\renewcommand{\familydefault}{\sfdefault}

\usepackage{amsmath,amssymb,mathtools}
\usepackage[version=4]{mhchem}
\usepackage{xcolor}

\definecolor{POSTECHRed}{RGB}{200,16,46}
\definecolor{POSTECHDark}{RGB}{25,25,25}

\setbeamercolor{normal text}{fg=black,bg=white}
\setbeamercolor{structure}{fg=POSTECHRed}
\setbeamercolor{alerted text}{fg=POSTECHRed}
\setbeamercolor{palette primary}{fg=white,bg=POSTECHRed}
\setbeamercolor{palette secondary}{fg=white,bg=POSTECHDark}
\setbeamercolor{palette tertiary}{fg=white,bg=POSTECHRed}
\setbeamercolor{palette quaternary}{fg=white,bg=POSTECHDark}
\setbeamercolor{frametitle}{fg=white,bg=POSTECHRed}
\setbeamercolor{title}{fg=white,bg=POSTECHRed}
\setbeamercolor{titlelike}{fg=white,bg=POSTECHRed}
\setbeamercolor{block title}{fg=white,bg=POSTECHRed}
\setbeamercolor{block body}{fg=black,bg=black!5}
\setbeamercolor{item projected}{fg=white,bg=POSTECHRed}
\setbeamercolor{subitem projected}{fg=white,bg=POSTECHRed}
\setbeamercolor{footer red}{fg=white,bg=POSTECHRed}
\setbeamercolor{footer dark}{fg=white,bg=POSTECHDark}

\setbeamertemplate{navigation symbols}{}
\setbeamertemplate{items}[circle]
\setbeamertemplate{blocks}[rounded][shadow=false]
\setlength{\parskip}{0.22em}
\setbeamerfont{itemize/enumerate body}{size=\large}
\setbeamerfont{itemize/enumerate subbody}{size=\normalsize}
\setbeamerfont{block body}{size=\normalsize}
\setbeamerfont{author in head/foot}{size=\scriptsize,series=\mdseries}
\setbeamerfont{title in head/foot}{size=\scriptsize,series=\mdseries}
\setbeamerfont{date in head/foot}{size=\scriptsize,series=\mdseries}
\setbeamerfont{title}{size=\LARGE,series=\bfseries}
\setbeamerfont{subtitle}{size=\large,series=\mdseries}

% Clean title page: author, institution, and date share one typographic level.
\setbeamertemplate{title page}{%
  \vbox{}
  \vfill
  \begin{centering}
    \begin{beamercolorbox}[sep=11pt,center,rounded=true]{title}
      {\usebeamerfont{title}\inserttitle\par}
      \vspace{0.32em}
      {\usebeamerfont{subtitle}\insertsubtitle\par}
    \end{beamercolorbox}
    \vspace{1.35em}
    {\normalsize\mdseries\insertauthor\par}
    \vspace{0.24em}
    {\normalsize\mdseries\insertinstitute\par}
    \vspace{0.24em}
    {\normalsize\mdseries\insertdate\par}
  \end{centering}
  \vfill
}

% Four-part footer requested for the simulation.
\setbeamertemplate{footline}{%
  \leavevmode%
  \hbox{%
    \begin{beamercolorbox}[wd=.31\paperwidth,ht=2.55ex,dp=1.05ex,leftskip=1.05em]{footer red}%
      \usebeamerfont{author in head/foot}20260918 Simulation
    \end{beamercolorbox}%
    \begin{beamercolorbox}[wd=.27\paperwidth,ht=2.55ex,dp=1.05ex,center]{footer dark}%
      \usebeamerfont{title in head/foot}Jungwoo Kim
    \end{beamercolorbox}%
    \begin{beamercolorbox}[wd=.20\paperwidth,ht=2.55ex,dp=1.05ex,center]{footer red}%
      \usebeamerfont{title in head/foot}POSTECH
    \end{beamercolorbox}%
    \begin{beamercolorbox}[wd=.22\paperwidth,ht=2.55ex,dp=1.05ex,rightskip=1.05em plus1fil]{footer dark}%
      \usebeamerfont{date in head/foot}\insertframenumber{} / \inserttotalframenumber
    \end{beamercolorbox}%
  }%
  \vskip0pt%
}

\newenvironment{problemframe}[2]{%
  \begin{frame}[t]{Problem #1 -- #2}%
  \large
}{%
  \end{frame}%
}

\newenvironment{solutionframe}[2]{%
  \begin{frame}[t]{Solution #1 -- #2}%
  \normalsize
}{%
  \end{frame}%
}

\newcommand{\answerbox}[1]{%
  \begin{block}{정답}
    \centering
    {\large #1}
  \end{block}
  \vspace{0.2em}%
}

\newenvironment{simchoices}{%
  \begin{enumerate}
    \setlength{\itemsep}{0.42em}%
    \setlength{\parskip}{0pt}%
}{%
  \end{enumerate}%
}

% MATH000-compatible notation used in this document.
\newcommand{\dd}{\mathop{}\!\mathrm{d}}
\newcommand{\N}{\mathbb{N}}
\newcommand{\Z}{\mathbb{Z}}
\newcommand{\Q}{\mathbb{Q}}
\newcommand{\R}{\mathbb{R}}
\newcommand{\C}{\mathbb{C}}
\DeclareMathOperator{\tr}{tr}
\DeclareMathOperator{\rk}{rk}
\DeclareMathOperator{\lcmop}{lcm}
